采用政策一致的 K=8 闭环方案，以 1 月预热后继承的 SOC 作为 2 月初状态。题设结果输出区间为 2025 年 2 月 1 日至 12 月 31 日，共 $334$ 日；该区间实际结算成本为 $15\,166\,538.46$ 元，其中普通计划购电成本为 $11\,883\,687.24$ 元，紧急购电成本为 $3\,282\,851.22$ 元。

上述数值体现的是锁定额度的计划成本与实际缺口的紧急成本之和，而非仅统计实际取用的普通电量。因而，风险储备过高会直接增大计划成本，储备不足又会通过 $5$ 倍价格的紧急购电体现出来；该成本结构正是模型选择预测结构和风险分位的依据。K=8 是在既定候选集合和闭环验证协议下保留的情景压缩规模，本文不将其表述为全局最优场景数。

为避免时间口径混淆，1 月预热期不并入上述题设输出成本。若需要报告 1--12 月完整核算，应另列预热期成本、每日 SOC 及 2 月 1 日初始 SOC，不得直接与本节的 $334$ 日结果相加或替换。
